\begin{center}
\section*{Abstract}
\addcontentsline{toc}{section}{Abstract}
\end{center}
This report involves three computational studies that integrate quantum chemical calculations and machine learning to solve important challenges in molecular chemistry and chemical data analysis. In the first experiment, the concept of aromaticity is analyzed by calculating resonance energies using thermodynamic electronic energies of different homocyclic and heterocyclic hydrocarbons and ions using Gaussian 16. Geometry optimizations and frequency analyses were performed using semi-empirical AM1 methods, followed by Hartree-Fock single point energy calculations with the 3-21G(*) basis set to assess aromatic stabilization.

The second experiment investigates the binding interactions of various Lewis acid-base adducts involving Boron $\&$ aluminum hydrides and halides as acid and nitrogen $\&$ phosphorus containing bases. Computational methods identical to those in the first study were employed, with modifications in basis sets to evaluate electronic and steric effects. Key parameters such as HOMO-LUMO energies, vertical electron attachment energy, reorganization energy, and proton affinity were analyzed to determine relative acid-base strengths and reactivity.

The third experiment applies machine learning techniques to analyze a physicochemical dataset of red and white wines quality dataset. Using Interactive Python and Jupyter notebooks, the study covers data pre-processing, visualization, and supervised learning models to derive patterns and build predictive frameworks. All three experiments show the utility of computational chemistry and data science in understanding molecular properties and handling experimental datasets.


